\section{A variance inequality for three rows}\label{row:section}

The estimate in this section concerns arbitrary binary matrices with three
rows and prescribed margins.  It will control resamplings on a vertex
triple, but does not require a graph interpretation.  The equivalent
conditional-projection inequality was proved by Fu, Qin and Wang
\cite[Theorem~5.1]{FQW26}.  We give a proof using label transpositions
and conditional variance, separating the randomness of the labels from
that of the block sizes. The count distribution and adjacent-total
coupling below correspond to their Lemmas~5.2--5.4; all steps are
included here.

The transposition and invariant-coordinate strategy for the fixed-block
argument was inspired by the symmetric-group heat-flow method of
Carlen, Lieb and Loss \cite[Section~2]{CLL06}.  The variance proof below
is self-contained; the three-row inequality and count coupling retain
the Fu--Qin--Wang attribution above.

\begin{lemma}[Fixed block sizes]\label{row:fixed-blocks}
Let $(U_1,U_2,U_3)$ and $(W_1,W_2,W_3)$ be independent uniform ordered
partitions of finite labeled sets $U$ and $W$, with all six block sizes
prescribed.  For arbitrary real functions $g_i=g_i(U_i,W_i)$,
\begin{equation}\label{row:fixed-variance}
 \operatorname{Var}\Bigl(\sum_{i=1}^3 g_i\Bigr)
 \le 2\sum_{i=1}^3\operatorname{Var}(g_i).
\end{equation}
\end{lemma}
\begin{proof}
Subtract the mean from each $g_i$.  On the finite conditional space, let
\[
 J=\sum_{\tau}(I-T_\tau),
\]
where $\tau$ ranges over transpositions of two labels within $U$ or
within $W$, and $T_\tau$ acts on functions by permuting those labels.
Every $T_\tau$ is a self-adjoint involution, so
\begin{equation}\label{row:transposition-energy}
 \langle f,Jf\rangle
 =\frac12\sum_\tau\mathbb E(f-T_\tau f)^2.
\end{equation}
The permutations of $U$ and $W$ act transitively on the ordered
partitions with their prescribed sizes.  Since transpositions generate
these permutations, the kernel of $J$ consists precisely of constants.
This also covers empty blocks and label sets; a one-state space has no
nonzero centered functions.

The space of functions of $(U_i,W_i)$ is invariant under $J$.
Indeed, a transposition within $U$ changes such a function only when it
exchanges a label in $U_i$ with one outside $U_i$.  The resulting block
is determined by $U_i$ and the two labels, regardless of which other
block contains the outside label.  The same observation applies to $W$.
Every spectral projection of the finite-dimensional operator $J$ is a
polynomial in $J$, and therefore preserves each row-function space.
Write $g_i=\sum_{\lambda>0}g_{i,\lambda}$ for its orthogonal spectral
decomposition.  For a fixed transposition, at most two of the three
block tuples change.  Hence, pointwise,
\[
 \left(\sum_i(g_{i,\lambda}-T_\tau g_{i,\lambda})\right)^2
 \le 2\sum_i(g_{i,\lambda}-T_\tau g_{i,\lambda})^2.
\]
Using \eqref{row:transposition-energy} gives
\[
 \lambda\left\|\sum_i g_{i,\lambda}\right\|^2
 \le 2\lambda\sum_i\|g_{i,\lambda}\|^2.
\]
Divide by $\lambda>0$ and sum over the orthogonal eigenspaces to obtain
\eqref{row:fixed-variance}.
\end{proof}

\begin{lemma}[Random block sizes]\label{row:counts}
Let $m,\alpha_1,\alpha_2,\alpha_3$ be nonnegative integers.  Give the
triples $x=(x_1,x_2,x_3)$ of nonnegative integers with
$x_1+x_2+x_3=m$ probabilities proportional to
\[
 \prod_{i=1}^3 w_i(x_i),\qquad
 w_i(t)=\frac1{t!(t+\alpha_i)!}.
\]
For arbitrary real functions $b_i$ of one coordinate,
\begin{equation}\label{row:count-variance}
 \operatorname{Var}\Bigl(\sum_{i=1}^3b_i(x_i)\Bigr)
 \le 2\sum_{i=1}^3\operatorname{Var}(b_i(x_i)).
\end{equation}
\end{lemma}
\begin{proof}
The case $m=0$ is immediate.  Suppose $m>0$.
We first compare the distributions of two coordinates with adjacent
totals, using the coupling of Fu, Qin and Wang \cite[Lemma~5.2]{FQW26}.
We include the argument.  Let $(Y_s,Z_s)$ have probabilities proportional to
$w_j(y)w_k(z)$ on $y+z=s$.  For $s\ge1$ and $0\le y\le s-1$, the
ratio of the probability of $Y_s=y$ to that of $Y_{s-1}=y$ is a positive
constant times
\[
 \frac{w_k(s-y)}{w_k(s-1-y)}
 =\frac1{(s-y)(s-y+\alpha_k)},
\]
which is nondecreasing in $y$.  Reweighting by a nondecreasing function
increases expectations of nondecreasing functions: for two such
functions $a,b$ and independent identically distributed $Y,Y'$, this
follows from
\[
 \operatorname{Cov}(a(Y),b(Y))
 =\frac12\mathbb E\bigl[(a(Y)-a(Y'))(b(Y)-b(Y'))\bigr]\ge0.
\]
Applying this observation on the common support, and then adding the
possible mass at $y=s$, proves $Y_{s-1}\le_{\mathrm{st}}Y_s$.
Interchanging the coordinates also gives
$Z_{s-1}\le_{\mathrm{st}}Z_s$, or equivalently
$Y_s\le_{\mathrm{st}}Y_{s-1}+1$.
Use the same uniform random variable in the inverse distribution
functions of $Y_{s-1}$ and $Y_s$.  Writing $Q$ for these quantile
functions, the identity $Q_{Y_{s-1}+1}(u)=Q_{Y_{s-1}}(u)+1$ shows that
both stochastic inequalities hold simultaneously:
\begin{equation}\label{row:covering-coupling}
 Y_{s-1}\le Y_s\le Y_{s-1}+1.
\end{equation}
Since these variables are integer-valued, setting $Z$ from its
respective total yields a coupling in which exactly one of the two
coordinates increases by one.

Now let $A$ map the direct sum of the three centered marginal function
spaces to the centered joint space by
$A(b_1,b_2,b_3)=\sum_i b_i(x_i)$.
The direct sum has squared norm $\sum_i\|b_i\|^2$.
To prove \eqref{row:count-variance}, it suffices to bound the largest
eigenvalue $\lambda$ of $A^*A$ by $2$.  A corresponding nonzero
centered eigenvector satisfies
\begin{equation}\label{row:count-eigenvector}
 \lambda b_i(t)
 =\mathbb E\left[\sum_{j=1}^3b_j(x_j)\,\middle|\,x_i=t\right],
 \qquad 0\le t\le m.
\end{equation}
Every value $0,\ldots,m$ has positive marginal probability.  Therefore
\[
 M=\max_{i,\,0\le t<m}|b_i(t+1)-b_i(t)|
\]
is positive: otherwise each component would be constant, and hence
zero by centering.

Conditioning on $x_i=t$ leaves the other two coordinates with total
$s=m-t$.  For adjacent values $t,t+1$, couple those coordinates using
\eqref{row:covering-coupling}.  Exactly one of them decreases by one,
so the change in their conditional sum of functions has absolute
value at most $M$.  Subtracting \eqref{row:count-eigenvector} at the two
values thus gives
\[
 |(\lambda-1)(b_i(t+1)-b_i(t))|\le M.
\]
Choose $i,t$ attaining $M$.  It follows that $|\lambda-1|\le1$, and
in particular $\lambda\le2$.
\end{proof}

\begin{lemma}[Three-row variance inequality]\label{row:variance}
Fix row and column sums for binary matrices with three rows, and take
the uniform distribution on the nonempty set of all matrices with
these prescribed margins.
Let $S_i$ denote row $i$.  For arbitrary real row functions $f_i$,
\begin{equation}\label{row:main-variance}
 \operatorname{Var}\Bigl(\sum_{i=1}^3f_i(S_i)\Bigr)
 \le 2\sum_{i=1}^3\operatorname{Var}(f_i(S_i)).
\end{equation}
\end{lemma}
\begin{proof}
Columns of sum $0$ or $3$ are deterministic; remove them and adjust
the row sums.  Let $U$ index the columns of sum $1$, assigned to the
row containing their unique one.  Let $W$ index the columns of sum $2$,
assigned to the row containing their unique zero.  Write $U_i,W_i$
for the resulting blocks.  The nondeterministic information in row
$i$ is exactly $(U_i,W_i)$.

Put $k_i=|W_i|$.  The adjusted row-sum constraints have the form
\[
 |U_i|=k_i+\delta_i
\]
for fixed integers $\delta_i$, since that row sum equals
$|U_i|+|W|-k_i$.  Conditional on the vector $k$, the assignments on $U$
and $W$ are independent uniform ordered partitions with prescribed
block sizes.  Lemma~\ref{row:fixed-blocks} therefore applies.

Counting these partitions shows that the law of $k$ has weights
proportional to
\[
 \prod_{i=1}^3\frac1{k_i!(k_i+\delta_i)!},
 \qquad k_i\ge0,\quad k_i+\delta_i\ge0,\quad \sum_i k_i=|W|.
\]
These conditions are sufficient for feasibility: the two sets of
nonnegative block sizes sum to $|W|$ and $|U|$, respectively.  Set
\[
 l_i=\max(0,-\delta_i),\qquad
 x_i=k_i-l_i,\qquad \alpha_i=|\delta_i|,\qquad
 m=|W|-\sum_i l_i.
\]
Nonemptiness implies $m\ge0$.  The support becomes all nonnegative
triples of total $m$, and the weights become
$\prod_i[x_i!(x_i+\alpha_i)!]^{-1}$.
Thus Lemma~\ref{row:counts} applies to functions of the individual
counts.

Conditional on $k$, the marginal blocks $U_i$ and $W_i$ are independent
uniform subsets of their respective sizes.  In particular
$\mathbb E[f_i(S_i)\mid k]$ depends only on $k_i$, or equivalently on
$x_i$.  Applying conditional variance, followed by
Lemmas~\ref{row:fixed-blocks} and~\ref{row:counts}, gives
\begin{align*}
 \operatorname{Var}\Bigl(\sum_i f_i(S_i)\Bigr)
 &=\mathbb E\operatorname{Var}\Bigl(\sum_i f_i(S_i)\,\Big|\,k\Bigr)
   +\operatorname{Var}\Bigl(\sum_i\mathbb E[f_i(S_i)\mid k]\Bigr)\\
 &\le 2\sum_i\mathbb E\operatorname{Var}(f_i(S_i)\mid k)
     +2\sum_i\operatorname{Var}(\mathbb E[f_i(S_i)\mid k])\\
 &=2\sum_i\operatorname{Var}(f_i(S_i)),
\end{align*}
as required.
\end{proof}
